% Author: Timofei Iuzhakov.
% Article section. Requires amsmath, amssymb, amsthm and the
% theorem, lemma, proposition, corollary, remark environments.
% Bibliography key: Lagarias. All labels in this section start with metric:.

\section{Digital rigidity and generic absence of three visits}
\label{metric:section}

Write
\[
 \Sigma=\left\{\sum_{j\geq0}e_j3^j:e_j\in\{0,1\}\right\}
 \subset\mathbb Z_3,\qquad G=1+3\mathbb Z_3,\qquad
 A=\Sigma\cap\mathbb Z_3^\times=1+3\Sigma.
\]
For a unit $b\in\mathbb Z_3^\times$, define
\[
 \begin{aligned}
 E_b^{(3)}&=\{\lambda\in\mathbb Z_3:
       b^n\lambda\in\Sigma\text{ for at least three distinct }n\geq0\},\\
 E_b&=\{\lambda\in\mathbb Z_3:
       b^n\lambda\in\Sigma\text{ for infinitely many }n\geq0\}.
 \end{aligned}
\]
All Haar measures on compact groups are normalized to have mass one.

\begin{theorem}\label{metric:main}
For Haar-almost every $b\in\mathbb Z_3^\times$,
\[
                         E_b^{(3)}=\{0\}.
\]
In particular, $E_b=\{0\}$ for Haar-almost every unit $b$.
The exceptional null set is independent of both the point $\lambda$ and
the three visit times.
\end{theorem}

This is a varying-base statement. The fixed-base exceptional sets were
studied by Lagarias \cite{Lagarias}; in particular, his
Theorem~1.6(iii) (Theorem~1.5(3) in the arXiv version~4) gives
$\dim_H E_2^{(3)}\geq (\log_3 2)/6$.
Thus the conclusion of Theorem~\ref{metric:main} cannot be imposed on the
prescribed base $2$. We do not attribute a varying-base conjecture to
Lagarias. Our proof uses a finite classification of extremal digital
planes, followed by an exact averaging identity over complete collections
of prefixes. Both ingredients are proved below.

\subsection{Full digital planes}

For $m\geq0$, put
\[
 D_m=\left\{\sum_{j<m}e_j3^j:e_j\in\{0,1\}\right\},\qquad
 B_m=\left\{\sum_{j<m}\sigma_j3^j:\sigma_j\in\{-1,1\}\right\},
 \qquad R_m=\frac{3^m-1}{2}.
\]
For unit coefficients $\alpha_1,\alpha_2,\alpha_3$ modulo $3^m$, let
\[
 P_m(C;\alpha)=\#\{(X_1,X_2,X_3)\in B_m^3:
               \alpha_1X_1+\alpha_2X_2+\alpha_3X_3\equiv C\pmod{3^m}\}.
\]
At most three of the eight next-digit triples satisfy a prescribed
residue modulo $3$. Consequently $P_m(C;\alpha)\leq3^m$.
We call a plane \emph{full} if equality holds. Equivalently, every
admissible node above the last level has three children.
Coordinate sign changes, permutations, and simultaneous multiplication
of the coefficients and target by a unit preserve the count.

\begin{theorem}[Finite digital rigidity]\label{metric:classification}
For $m\geq1$, a plane is full if and only if, after the operations just
listed, its coefficients and target have one of the following forms
modulo $3^m$:
\begin{align*}
 &(1,1,1),&& C\in B_m;\tag{a}\label{metric:classa}\\
 &(1,1,1+\tau3^g),&& C\in B_g,
       \quad 1\leq g<m,\quad\tau\in\{-1,1\}.\tag{b}\label{metric:classb}
\end{align*}
In (b), $B_g$ means its $2^g$ ordinary integer representatives reduced
modulo $3^m$, not all lifts of $B_g$ from modulus $3^g$.
\end{theorem}

We first establish the coefficient coincidence that underlies the
classification.

\begin{lemma}\label{metric:coincidence}
Normalize the coefficients by signs so that all three are $1$ modulo
$3$. If a depth-$m$ plane is full, two of these normalized coefficients
are equal modulo $3^m$.
\end{lemma}

\begin{proof}
Divide by the first coefficient. Use binary coordinates through the
bijection $X=2x-R_m$ from $D_m$ to $B_m$; modulo $3^m$ this is
$X\equiv2x+1/2$. For fixed coefficients let $F_l$ be the full binary
targets at depth $l$. Reduction sends $F_l$ into $F_{l-1}$, and each
of the three admitted first-digit choices below a full target has a
child target in $F_{l-1}$.

If all coefficients equal $1$ modulo $3^m$, there is nothing to prove.
Otherwise let $k<m$ be their first unequal depth, so that
\[
 \alpha_i=1+3^k a_i,\qquad a_1=0,
 \qquad a_1,a_2,a_3\text{ are not all equal modulo }3.
\]
Here $k\geq1$. For three equal binary coefficients, a target has three
first children precisely when its first digit is $1$ or $2$; all three
children have the same residual target. Induction gives
\[
                         F_k=R_k+D_k.
\]
For the canonical representative $V_0\in F_k$, every admitted triple
of depth-$k$ prefixes satisfies $x_1+x_2+x_3=V_0$ as an ordinary
equality: the sum is at most $3R_k<V_0+3^k$, and a negative wraparound
is impossible since $0\leq V_0<3^k$.

Consider a lift $V_0+3^k v$. Its next required residue, for an old
prefix $(x_1,x_2,x_3)$, is
$v-\sum_i a_i(x_i\bmod3)$. This residue must be nonzero.
If $V_0\equiv1\pmod3$, the three first-digit choices exclude
$v\in\{a_1,a_2,a_3\}$; if $V_0\equiv2\pmod3$, they exclude
$v\in\{a_1+a_2+a_3-a_j:1\leq j\leq3\}$.
All three choices occur in the full prefix tree. Three distinct
$a_i$ residues therefore permit no full lift. A split into two equal
residues and one other permits exactly one lift. Since a full
depth-$m$ target exists, the latter case applies, and $F_{k+1}$ is
injective under reduction modulo $3^k$. Fix the equal pair $i,j$;
then $\alpha_i\equiv\alpha_j\pmod{3^{k+1}}$.

This target injectivity persists at every depth $l\geq k+1$.
Indeed, suppose it holds at $l-1$, and take $V,V'\in F_l$ with
the same reduction modulo $3^k$. Their reductions modulo
$3^{l-1}$ agree. Choose the same admitted first-digit triple for
both targets. The child targets differ by $(V-V')/3$, which is zero
modulo $3^{l-2}$, and hence modulo $3^k$ when $l\geq k+2$.
Injectivity in $F_{l-1}$ gives $V\equiv V'\pmod{3^l}$.

If $m=k+1$, the selected coefficients already coincide modulo $3^m$.
Otherwise, among the three first children of a full depth-$m$ target
there are two whose targets differ by
$\pm(\alpha_i-\alpha_j)/3$. They agree modulo $3^k$.
Injectivity in $F_{m-1}$ forces their equality modulo $3^{m-1}$,
and therefore $\alpha_i\equiv\alpha_j\pmod{3^m}$.
\end{proof}

\begin{proof}[Proof of Theorem~\ref{metric:classification}]
By Lemma~\ref{metric:coincidence}, we may work with balanced
coefficients $(1,1,t)$, where $t\equiv1\pmod3$.
If $t\equiv1\pmod{3^m}$, every admitted child map is
$C\mapsto(C-\sigma)/3$, where $\sigma\in\{-1,1\}$ represents
$C\bmod3$. The plane is full exactly when $C\in B_m$.

Otherwise set $g=v_3(t-1)<m$, using any unit lift of $t$ to
$\mathbb Z_3$. The child maps of a full state are
\begin{equation}\label{metric:children}
 C\longmapsto\frac{C-\sigma t}{3}\quad\text{(twice)},
 \qquad
 C\longmapsto\frac{C-\sigma(2-t)}{3}\quad\text{(once)}.
\end{equation}
Put
\[
 \delta=\frac{2(1-t)}t,\qquad
 d\in\{-1,1\},\quad d\equiv\delta/3^g\pmod3,
 \qquad\epsilon=-d,\qquad D=1-\epsilon3^g.
\]
For any full state of depth $l$, always following the first map in
\eqref{metric:children} shows that $z=C/t$ has balanced digits
$\sigma_0,\ldots,\sigma_{l-1}$. Change only the branch at time $j$
to the second map and then follow first maps again. The signs through
time $j$ are unchanged. Telescoping the digit expansions gives
\[
                     z-\delta\sigma_j3^j\in B_l\pmod{3^l}.
\]
When $j+g<l$, the first affected digit is at $j+g$. It cannot be zero
modulo $3$, so
\begin{equation}\label{metric:period}
                         \sigma_{j+g}=\epsilon\sigma_j.
\end{equation}
It follows, at every finite depth $l\geq g$, that
\begin{equation}\label{metric:periodicform}
                  z\equiv P/D\pmod{3^l}\quad\text{for some }P\in B_g.
\end{equation}
At $l=g$ this only records the first $g$ digits, since $D\equiv1$
modulo $3^g$; at greater depths it is precisely
\eqref{metric:period}. No infinite extension has been assumed.

Apply \eqref{metric:periodicform} to the two different children of
the depth-$m$ root. Their normalized targets satisfy
\[
 z_0\equiv P_0/D,\qquad z_1\equiv P_1/D\pmod{3^{m-1}},
 \qquad P_0,P_1\in B_g,
\]
and, writing $\sigma=\sigma_0$,
\begin{equation}\label{metric:childdiff}
               P_1-P_0\equiv-\sigma\delta D/3\pmod{3^{m-1}}.
\end{equation}
This difference has valuation $g-1$. Thus the two balanced words
have equal lower $g-1$ digits and opposite highest digits.
The highest digit of $P_0$ is the root digit $\sigma_g$, which
equals $\epsilon\sigma$ by \eqref{metric:period}. Therefore, as
an ordinary integer equality,
\[
                    P_1-P_0=-2\epsilon\sigma3^{g-1}.
\]
Equation~\eqref{metric:childdiff} yields
$\delta D\equiv2\epsilon3^g\pmod{3^m}$. The identity
\[
                     \delta D-2\epsilon3^g=2(D-t)/t
\]
then gives $t\equiv D\pmod{3^m}$, and
\eqref{metric:periodicform} at the root gives $C\in B_g$ at the
full modulus $3^m$. This also proves the boundary case $m=g+1$:
the two child words have length $g$, and their differing digit is
still at position $g-1$.

Conversely, take exact representatives $t=1-\epsilon3^g$ and
$C\in B_g$. The maps \eqref{metric:children} become
\[
           \frac{C-\sigma}{3}\mathbin{\pm}\epsilon\sigma3^{g-1}.
\]
Both remain in $B_g$: they shift its lower $g-1$ remaining digits
and insert an allowed highest sign. Every reachable state is a unit
and has three children. The plane is full at all depths, and reduction
gives (b). The common-coefficient case was already proved.
\end{proof}

For later use define the two-input count
\begin{equation}\label{metric:Q}
 \begin{aligned}
 Q_m(U;a,b)&=\#\{(x,y)\in D_m^2:
                     U+ax+by\bmod3^m\in D_m\},\\
 H_m(a,b)&=\max_{U\bmod3^m}Q_m(U;a,b),
 \end{aligned}
\end{equation}
where $a,b$ are units. The balanced-coordinate conversion gives
\begin{equation}\label{metric:conversion}
 Q_m(U;a,b)=P_m\bigl((a+b-1)/2-2U;(a,b,-1)\bigr).
\end{equation}
In particular, $H_m(a,b)\leq3^m$.

\begin{corollary}\label{metric:inverse}
If $a,b\in\mathbb Z_3^\times$ and $H_m(a,b)=3^m$ for every
$m\geq1$, then
\[
                   a\in\{1,-1\},\quad\text{or}\quad
                   b\in\{1,-1\},\quad\text{or}\quad a\in\{b,-b\}.
\]
\end{corollary}
\begin{proof}
At every depth Lemma~\ref{metric:coincidence} gives one of these six
equalities modulo $3^m$. One equality occurs at unbounded depths,
and hence holds in $\mathbb Z_3$.
\end{proof}

\begin{corollary}[Enumeration and a rational height bound]
\label{metric:effective}
There are exactly $24m-20$ ordered unit pairs $(a,b)$ modulo $3^m$
for which $H_m(a,b)=3^m$, and exactly $28\cdot2^m-48$ full triples
$(a,b,U)$.
Moreover, let $a=A/D$, $b=B/D$, where $A,B,D$ are ordinary integers
prime to $3$ and $\max(|A|,|B|,|D|)\leq H$. If the pair is not full
at every depth, then
\[
                   3^m>4H^2\quad\Longrightarrow\quad
                   H_m(a,b)\leq3^m-1.
\]
\end{corollary}
\begin{proof}
Normalize $a,b$ separately by signs to be $1$ modulo $3$.
Theorem~\ref{metric:classification} and \eqref{metric:conversion}
give the normalized pairs
\[
 (1,1),\qquad (t,1),\ (1,t),\ (t^{-1},t^{-1}),
 \qquad t=1\pm3^g,\quad1\leq g<m.
\]
They are distinct: each nonconstant triple has a unique coincident
pair, and $v_3(t-1)=g$ separates the levels. There are $6m-5$
normalized pairs and four sign choices. The common pair has $2^m$
targets per sign choice; each nonconstant pair of level $g$ has
$2^g$. Hence the target count is
$4\cdot2^m+24\sum_{g=1}^{m-1}2^g=28\cdot2^m-48$.

Suppose now $3^m>4H^2$ and the rational pair is full at depth $m$.
Finite coefficient coincidence makes one of $A\pm B$, $A\pm D$,
$B\pm D$ divisible by $3^m$. Its absolute value is at most $2H$,
so it is zero. After normalization the triple is $(1,1,T/S)$
with $|S|,|T|\leq H$. If $T=S$, the game is full at all depths.
Otherwise put $g=v_3(T-S)$; then $3^g\leq2H<3^m$.
Classification gives $T-S(1\pm3^g)\equiv0\pmod{3^m}$, while
\[
                 |T-S(1\pm3^g)|\leq2H+2H^2\leq4H^2<3^m.
\]
This is again an exact equality, placing the coefficients in an
all-depth full family. The assumed nonfullness proves the implication.
\end{proof}

\subsection{Complete prefix cubes and Haar measure}

We use the elementary power operation $\rho^N$, for $\rho\in G$ and
$N\in\mathbb Z_3$. It is defined by the convergent binomial series,
or equivalently by $\exp(N\log\rho)$. The series for logarithm and
exponential give inverse group isomorphisms between $G$ and
$3\mathbb Z_3$. In particular a unit power is an automorphism of $G$.
The coefficients $\binom Nk$ are integral: approximate $N$ by
nonnegative integers in $\mathbb Z_3$ and use continuity of each
binomial polynomial. These series also show that
$\rho\mapsto\rho^N$ is $1$-Lipschitz. We will use these facts only
on principal units.

For fixed $N\in\mathbb Z_3\setminus\{0,1\}$ define
\[
 \mathcal B_N=\{\rho\in G:\text{some }\mu\in A\text{ satisfies }
                           \rho\mu,\rho^N\mu\in A\}.
\]
This set is compact, being the projection of a closed subset of
$G\times A$.

\begin{proposition}\label{metric:fixedratio}
For every fixed $N\in\mathbb Z_3\setminus\{0,1\}$, the set
$\mathcal B_N$ has Haar measure zero in $G$.
\end{proposition}

\begin{proof}
If $N\equiv2\pmod3$, then $\mathcal B_N\subseteq\{1\}$.
Indeed, write $\rho=1+3^g t\ne1$ with $g\geq1$ and $t$ a unit.
The three numbers $\mu,\rho\mu,\rho^N\mu$ have equal digits below
$g$, whereas their digits at $g$ differ by $0,t,2t$ modulo $3$
because $\mu\equiv1\pmod3$. One of these three digits is $2$.

Next suppose $N=3^s r$ with $s\geq1$ and $r$ a unit. Set
\begin{equation}\label{metric:coefficients}
 a(\rho)=(1-N)\rho^N,\qquad b(\rho)=r\rho^{N-1},\qquad
 \Phi(\mu,\nu)=\mu^{1-N}\nu^N.
\end{equation}
Both coefficients are units. A level-$q$ node is a pair
$(\mu,\nu)\in A_q^2$, where $A_q=1+3D_{q-1}$, for which
$\Phi(\mu,\nu)\bmod3^q\in A_q$. Representatives in $A_q$ are
ordinary nonnegative integers. Its ratio is $\rho=\nu/\mu\in G$.

Fix $m\geq1$ and $q\geq m+s+1$. Group pairs by fixing $\mu$ and
the digits of $\nu$ below position $q-s$. Write the resulting pairs as
\[
            (\mu,\nu_e),\qquad
            \nu_e=\nu_0+3^{q-s}e,\qquad e\in D_s.
\]
Every group containing a node contains all its $2^s$ possible nodes.
More precisely, with $\rho_0=\nu_0/\mu$,
\begin{equation}\label{metric:cubeshift}
 \Phi(\mu,\nu_e)\equiv\Phi(\mu,\nu_0)
               +3^q b(\rho_0)e\pmod{3^{q+m+1}}.
\end{equation}
To prove this, expand in $3^{q-s}e/\nu_0$. From
$\binom Nk=(N/k)\binom{N-1}{k-1}$ we obtain
$v_3\binom Nk\geq s-v_3(k)$. For $k\geq2$ the valuation of the
$k$th term is at least
\[
 k(q-s)+s-v_3(k)\geq2q-s\geq q+m+1;
\]
the first inequality follows from $v_3(k)\leq k-2$ and $q-s\geq1$.
This proves \eqref{metric:cubeshift}. In addition, all ratios in the
group agree modulo $3^{m+1}$.

Let $U$ be the block of $m$ digits of $\Phi(\mu,\nu_0)$ beginning
at position $q$. The descendants of a group have the form
\[
                     (\mu+3^q x,\nu_e+3^q y),\qquad x,y\in D_m.
\]
Integral binomial coefficients and $2q\geq q+m+1$ give their exact
new-block condition
\begin{equation}\label{metric:cubelinear}
              U+be+ax+3^sby\bmod3^m\in D_m,
\end{equation}
where $a,b$ are evaluated at $\rho_0$. They can be used throughout
the group by the Lipschitz property and the common ratio residue.
The lower $q$ digits remain admissible.

As $e\in D_s$ and $y\in D_m$ vary, $z=e+3^s y$ runs bijectively
over $D_{m+s}$. Each residue in $D_m$ has exactly $2^s$ such lifts
modulo $3^m$, also when $m<s$. Summing \eqref{metric:cubelinear}
therefore gives the exact descendant count
\begin{equation}\label{metric:cubecount}
                \#\{\text{descendants of the group}\}
                         =2^s Q_m(U;a,b).
\end{equation}

Fix a ratio class $\mathcal C\subset G$ modulo $3^{m+1}$.
The coefficients are constant modulo $3^m$ on this class. Write
$H_{\mathcal C}=H_m(a,b)$, and let $T_q(\mathcal C)$ count its
level-$q$ nodes. All the preceding groups stay in $\mathcal C$,
as do their descendants. Since each group has $2^s$ original nodes,
\eqref{metric:cubecount} implies
\begin{equation}\label{metric:contraction}
             T_{q+m}(\mathcal C)\leq H_{\mathcal C}T_q(\mathcal C),
             \qquad q\geq m+s+1.
\end{equation}
If $H_{\mathcal C}<3^m$, iterate with $q_j=q_0+jm$ to obtain
$T_{q_j}(\mathcal C)\leq T_{q_0}(\mathcal C)H_{\mathcal C}^j$.

For a ratio $\rho$ modulo $3^q$, let $n_q(\rho)$ count the
$\mu\in A_q$ for which $\rho\mu$ and $\rho^N\mu$ lie in $A_q$
modulo $3^q$. The substitution $\nu\equiv\rho\mu\pmod{3^q}$
is a bijection, and $G$ has $3^{q-1}$ residue classes modulo $3^q$.
Consequently
\[
 \int_{\mathcal C}n_q(\rho)\,d\rho
                     =\frac{T_q(\mathcal C)}{3^{q-1}}.
\]
Every $\rho\in\mathcal B_N$ has $n_q(\rho)\geq1$ at each depth.
It follows that
\[
 \operatorname{Haar}(\mathcal B_N\cap\mathcal C)
 \leq\frac{T_{q_j}(\mathcal C)}{3^{q_j-1}}\longrightarrow0
 \qquad(H_{\mathcal C}<3^m).
\]
Taking the countable union over finite windows and ratio classes
reduces the proof to the residual set
\[
 \mathcal F_N=\{\rho\in G:H_m(a(\rho),b(\rho))=3^m
                                      \text{ for every }m\geq1\}.
\]

In fact $\mathcal F_N$ contains at most two points. Let
$\sigma\in\{-1,1\}$ represent $r\bmod3$. Since
$\rho^N\equiv1\pmod{3^{s+1}}$,
\[
                    a(\rho)-1\equiv-N\pmod{3^{s+1}},
                    \qquad v_3(a(\rho)-1)=s.
\]
Thus $a\ne1$, and $a\equiv1\pmod3$ also excludes $a=-1$.
Corollary~\ref{metric:inverse} leaves only $b=\sigma$ or
$a=\sigma b$. The first has the unique solution
$\rho_1=(\sigma/r)^{1/(N-1)}\in G$, because $N-1$ is a unit.
The second, using $a/b=((1-N)/r)\rho$, has the unique solution
$\rho_2=\sigma r/(1-N)\in G$. Hence
$\mathcal F_N\subseteq\{\rho_1,\rho_2\}$, proving nullity.
No assertion that either candidate is full is needed.

Finally, if $N\equiv1\pmod3$, exchange $\mu$ and $\nu=\rho\mu$.
The ratio becomes $\rho^{-1}$ and the exponent becomes $1-N$,
which is divisible by $3$ and different from $0,1$.
Thus inversion maps $\mathcal B_N$ into $\mathcal B_{1-N}$.
Inversion preserves Haar measure, and the preceding case applies.
\end{proof}

\subsection{One null set for all points and times}

\begin{proof}[Proof of Theorem~\ref{metric:main}]
For every nonzero integer $e$, the map $P_e:G\to G$ given by
$P_e(\beta)=\beta^e$ has null preimages of Haar-null sets.
Indeed, in logarithmic coordinates it is $x\mapsto ex$ on
$3\mathbb Z_3$, an injective homomorphism onto an open subgroup.
Its pushforward measure is normalized Haar measure on that subgroup,
and is therefore absolutely continuous with respect to Haar measure
on $G$.

By Proposition~\ref{metric:fixedratio}, the set
\[
 \mathcal N=
 \bigcup_{e\in\mathbb Z\setminus\{0\}}
 \ \bigcup_{N\in(\mathbb Q\cap\mathbb Z_3)\setminus\{0,1\}}
                         P_e^{-1}(\mathcal B_N)
\]
is Haar-null. Fix $\beta\in G\setminus\mathcal N$. If
$\lambda\ne0$ had three distinct visit times $n_0,n_1,n_2$, put
$v=v_3(\lambda)$. Interchange $n_1,n_2$ if necessary so that
\[
 v_3(e)\leq v_3(e'),\qquad e=n_1-n_0,\qquad e'=n_2-n_0.
\]
Then $N=e'/e\in(\mathbb Q\cap\mathbb Z_3)\setminus\{0,1\}$,
and the quantities
\[
             \mu=3^{-v}\beta^{n_0}\lambda,\qquad\rho=\beta^e
\]
satisfy $\mu,\rho\mu,\rho^N\mu\in A$.
The identity $\rho^N=\beta^{e'}$ follows from logarithms.
Thus $\rho\in\mathcal B_N$, contradicting the choice of $\beta$.
This proves the theorem on $G$ simultaneously for all nonzero
$\lambda$ and all triples of times.

For $b\equiv-1\pmod3$, the visit times of a fixed nonzero point
have a common parity: the leading unit digit of $b^n\lambda$ must
be $1$. Three visits therefore give three visits of $b^{r_0}\lambda$
under the base $\beta=b^2\in G$, for some $r_0\in\{0,1\}$.
Squaring maps the coset $-1+3\mathbb Z_3$ bijectively onto $G$
and has null preimages of null sets; write $b=-u$ and use the unit
power automorphism $u\mapsto u^2$. This proves the statement on
all units. Finally, $0$ always visits $\Sigma$ and
$E_b\subseteq E_b^{(3)}$.
\end{proof}

\begin{remark}\label{metric:scope}
The averaging in \eqref{metric:cubecount} is over a complete group
of prefixes. It does not give a contraction on a prescribed ratio
fiber. Accordingly, neither Theorem~\ref{metric:main} nor the
effective finite-window bound in Corollary~\ref{metric:effective}
asserts finiteness of visits for a specified base and a specified
point. The proofs above are independent of finite enumeration;
the accompanying exact checks verify small instances of the
classification and the complete-cube identity.
\end{remark}
